\pagina{Conclusiones}
Se implementaron las seis familias requeridas con un contrato común de ejecución y una representación que preserva permutaciones. La función de costo cuenta aristas faltantes, incluido el cierre; el éxito requiere un ciclo certificado y no solo una ruta bien formada o un fitness favorable.

La campaña histórica completó su protocolo y expuso límites de las variantes iniciales: cero éxitos en 900 primarios de evaluación y escala. Las mejoras posteriores y su evaluación independiente produjeron \TotalSuccesses{} éxitos en \TotalAttempts{} intentos y \SolvedGraphs{} grafos distintos certificados. Son etapas de condiciones diferentes que se conservan separadas, sin adjudicar la mejora a un único cambio.

ABC fue la familia que mejor se adaptó al corte observado por cobertura y éxitos: alcanzó los 15 grafos resueltos, con 116 éxitos en 1\,440 intentos sobre cobertura completa. Esta conclusión se restringe a las instancias, parámetros y recursos evaluados. Los tiempos cercanos al plazo, los equipos distintos y los agotamientos impiden sostener una superioridad temporal universal.

Las exploraciones complementarias certificaron graph4 y graph35 con la variante guiada y graph1, graph39, graph43, graph47 y graph52 con reparación v4. La evaluación v4 quedó cerrada en \RepairAttempts{} intentos, con \RepairSuccesses{} éxitos y \RepairErrors{} interrupciones, sobre \RepairComplete{} grafos completos y \RepairPartial{} parciales. ABC, GA y AS aportaron los ciclos de ese protocolo. La reparación produjo soluciones válidas en esta selección, sin que el diseño permita aislar su efecto causal.

El contraste FHCP se apoyó en fuentes primarias, distinguiendo caminos de ciclos y señalando presupuestos y condiciones ausentes. No se presentaron las referencias oficiales ni las publicaciones como ejecuciones propias, ni se calcularon aceleraciones entre entornos no equivalentes.

La unión de las entregas contiene \UnionCertifiedGraphs{} grafos distintos certificados, por encima de la referencia aproximada de 20 y por debajo de la meta original de 60. La evaluación distribuida conserva sus cinco parciales, 12 incidencias y 51\,254 combinaciones fuera de alcance; v4 conserva otras \RepairOmitted{} combinaciones no ejecutadas. Estos conteos pertenecen a planes diferentes y no se suman como cobertura homogénea. Se cierra la entrega en los cortes documentados, sin retomar las corridas omitidas. El principal aporte es una comparación verificable de implementaciones reales, con datos recuperables y tablas/figuras regenerables, además de conclusiones ajustadas al alcance de la evidencia.
